
Il suffit donc, pour achever la démonstration, de noter que :
@@M0@@

Pour montrer que @@M1@@ est nul, on prend @@M2@@, et @@M3@@ il existe alors @@M4@@ tel que @@M5@@, et @@M6@@; comme @@M7@@ est de dimension @@M8@@, cela entraine @@M9@@, donc @@M10@@ est acyclique, et @@M11@@.

\textbf{Lemme 6.7.} \uline{Soit} @@M12@@ \uline{un schéma affine noethérien}, @@M13@@ \uline{un nombre fini de points de} @@M14@@, @@M15@@ \uline{un} @@M16@@\uline{-Module inversible}. \uline{Alors}, \uline{il existe une section} @@M17@@ \uline{telle que pour tout} @@M18@@, @@M19@@ \uline{soit une base de} @@M20@@.

Posons @@M21@@. Pour tout @@M22@@, soit @@M23@@ un point fermé de @@M24@@ tel que @@M25@@. Les @@M26@@ correspondent à des idéaux maximaux @@M27@@ de @@M28@@, qu'on peut supposer tous distincts, quitte à changer l'ensemble d'indices de façon que chacun des @@M29@@ intervienne une fois et une seule dans la famille @@M30@@. Comme @@M31@@ est inversible, on a pour tout @@M32@@ un isomorphisme @@M33@@ d'où on déduit un isomorphisme @@M34@@, puis @@M35@@, soit encore @@M36@@. En relevant ce dernier isomorphisme, on obtient un homomorphisme @@M37@@ tel que l'homomorphisme @@M38@@ correspondant soit pour tout @@M39@@ un isomorphisme modulo @@M40@@, donc un isomorphisme d'après Nakayama. L'image de @@M41@@ par cet homomorphisme donne donc une section @@M42@@ de @@M43@@ qui est une base de @@M44@@ pour tout @@M45@@. Comme l'ensemble des points où @@M46@@ est une base de @@M47@@ est un ouvert de @@M48@@, @@M49@@ est aussi une base de @@M50@@ pour @@M51@@.

\textbf{Lemme 6.8.} \uline{Soient} @@M52@@ \uline{un schéma noethérien ayant un faisceau ample} @@M53@@, @@M54@@ \uline{un} @@M55@@\uline{-Module inversible}, @@M56@@ \uline{des points de} @@M57@@. \uline{Alors il existe un}
